\subsection{El Problema de la Mochila Clásico}

El problema de la mochila 0/1 es un problema de optimización combinatoria NP-difícil
ampliamente estudiado en la literatura. Dado un conjunto $\mathcal{N} = \{1, 2, \ldots, n\}$
de $n$ ítems, cada uno con un valor $v_i$ y un peso $w_i$, y una mochila de capacidad
máxima de peso $W$, el problema consiste en seleccionar un subconjunto de ítems que
maximice el valor total sin superar la capacidad:

\begin{equation}
  \text{Maximizar} \quad Z = \sum_{i=1}^{n} v_i x_i
  \label{eq:obj_clasico}
\end{equation}
\begin{equation}
  \text{sujeto a} \quad \sum_{i=1}^{n} w_i x_i \leq W, \quad x_i \in \{0,1\}
  \label{eq:cap_clasico}
\end{equation}

\subsection{Versión Extendida Implementada}

Para esta actividad se implementa una versión extendida donde cada ítem $i$ posee:
valor $v_i$, peso $w_i$, volumen $u_i$ y categoría $c_i$. La mochila añade las siguientes
restricciones sobre las restricciones clásicas:

\begin{enumerate}[label=\textbf{R\arabic*.}]
  \item \textbf{Restricción de peso}: $\displaystyle\sum_{i=1}^{n} w_i x_i \leq W$
  \item \textbf{Restricción de volumen}: $\displaystyle\sum_{i=1}^{n} u_i x_i \leq V$
  \item \textbf{Restricciones por categoría}: para cada categoría $k$,
        $\text{min}_k \leq \displaystyle\sum_{i : c_i = k} x_i \leq \text{max}_k$
  \item \textbf{Incompatibilidades}: para ciertos pares $(a,b)$, $x_a + x_b \leq 1$
  \item \textbf{Dependencias}: si $x_a = 1$, entonces $x_b = 1$
\end{enumerate}

La formulación extendida queda:

\begin{equation}
  \text{Maximizar} \quad Z = \sum_{i=1}^{n} v_i x_i
\end{equation}
\begin{align}
  \text{sujeto a} \quad &\sum_{i} w_i x_i \leq W \notag\\
                        &\sum_{i} u_i x_i \leq V \notag\\
                        &\text{min}_k \leq \sum_{i: c_i=k} x_i \leq \text{max}_k \quad \forall k \notag\\
                        &x_a + x_b \leq 1 \quad \forall (a,b) \in \mathcal{I} \notag\\
                        &x_a \leq x_b \quad \forall (a,b) \in \mathcal{D} \notag\\
                        &x_i \in \{0,1\} \notag
\end{align}

donde $\mathcal{I}$ es el conjunto de pares incompatibles y $\mathcal{D}$ el conjunto de
pares con dependencia.

\subsection{Generación de Instancias}

Las instancias se generan con tres tamaños para cubrir distintos escenarios de rendimiento:

\begin{table}[H]
  \centering
  \caption{Configuración de las tres instancias de prueba}
  \label{tab:instancias}
  \begin{tabular}{lccccc}
    \toprule
    \textbf{Instancia} & \textbf{Ítems} & \textbf{$W$} & \textbf{$V$} &
    \textbf{Incomp.} & \textbf{Dep.} \\
    \midrule
    Pequeña (\texttt{small})  & 100    &
      $0.40 \times \sum w_i$ & $0.40 \times \sum u_i$ & $\approx50$  & $\approx30$  \\
    Mediana (\texttt{medium}) & 1,000  &
      $0.40 \times \sum w_i$ & $0.40 \times \sum u_i$ & $\approx500$ & $\approx300$ \\
    Grande (\texttt{large})   & 10,000 &
      $0.40 \times \sum w_i$ & $0.40 \times \sum u_i$ & $\approx5000$& $\approx3000$\\
    \bottomrule
  \end{tabular}
\end{table}

Las capacidades máximas se fijaron al 40\,\% de la suma total de pesos y volúmenes
respectivamente, asegurando que existan soluciones factibles no triviales (instancias de
dificultad media, según la recomendación de la actividad).